\begin{theorem}[Nonzero regular states with inserted bond operators]
    \label{thm:peps_regular_twisted_state_nonzero}
    \lean{TNLean.PEPS.regularProjectorTwistedRegionMatrix_one_column_ne_zero,
        TNLean.PEPS.openRegionWeight_regularTwistedSite_ne_zero_of_isGInjective,
        TNLean.PEPS.stateCoeff_regularTwistedSite_ne_zero_of_isGInjective,
        TNLean.PEPS.IsGIsometric.stateCoeff_torusRegularTwistedSite_ne_zero}
    \leanok
    \uses{thm:peps_regular_twisted_region_accessibility,
        lem:peps_regular_torus_incident_isometry}
    Let $G$ be finite, and let the local tensors on a finite graph be
    regular $G$-injective. Insert arbitrary group-valued operators on
    the oriented bonds. The resulting closed PEPS vector is nonzero.
    More generally, each open-region vector with identity boundary
    labels is nonzero, as is the corresponding canonical projector column.
    In particular, on a torus with both periods at least three, every
    actual regular $G$-isometric state with inserted bond operators is nonzero.
    This is an auxiliary nonvanishing result for the flux measurements
    of \cite[Theorem~6.15]{Schuch2010PEPS}, without a parent-Hamiltonian claim.
\end{theorem}
\begin{proof}\leanok
    Set all bond labels and vertex translations to the identity, and
    choose each canonical physical row to be the resulting inserted
    half-edge labels. This gives a positive term in the actual
    canonical coefficient sum. All remaining terms are nonnegative.
    The derived local physical inverse sends the original contraction
    to this canonical column, so the original vector cannot vanish.
    For the full region, the exterior multiplicity is one and the
    open-region contraction is the closed PEPS vector. The native
    torus assertion follows from the incident-leg expression of local
    regular isometry.
\end{proof}
